\section*{Capítulo \ref{ch:b1:functions} --- Funções Usuais}

\begin{solution}{exo:b1:functions:1}
$\arcsin\frac{\sqrt 3}{2} = \frac\pi3$;
$\;\arccos\bigl(-\frac12\bigr) = \frac{2\pi}{3}$;
$\;\arctan(-1) = -\frac\pi4$.

$\arcsin\bigl(\sin\frac{5\pi}{6}\bigr)$: $\sin\frac{5\pi}{6} =
\frac12$, e o ângulo de $\intcc{-\frac\pi2}{\frac\pi2}$ cujo seno vale
$\frac12$ é $\frac\pi6$ (e não $\frac{5\pi}{6}$).

$\arccos\bigl(\cos\frac{7\pi}{4}\bigr)$: $\cos\frac{7\pi}{4} =
\frac{\sqrt 2}{2}$, e o ângulo de $\intcc{0}{\pi}$ com esse cosseno
é $\frac\pi4$.
\end{solution}

\begin{solution}{exo:b1:functions:2}
$f(x) = \arcsin(2x - 1)$ exige $-1 \leq 2x - 1 \leq 1$, isto é, $x
\in \intcc{0}{1}$. Em $\intoo{0}{1}$, a regra da cadeia e a
\cref{prop:b1:functions:arcderiv} dão
\[
f'(x) = \frac{2}{\sqrt{1 - (2x-1)^2}} = \frac{2}{\sqrt{4x - 4x^2}}
= \frac{1}{\sqrt{x(1 - x)}} .
\]
$g(x) = \arctan\sqrt x$ está definida para $x \geq 0$ e, para $x > 0$:
\[
g'(x) = \frac{1}{1 + x} \cdot \frac{1}{2\sqrt x}
= \frac{1}{2\sqrt x\,(1 + x)} .
\]
\end{solution}

\begin{solution}{exo:b1:functions:3}
$\cosh x + \sinh x = \frac{\eu^x + \eu^{-x}}{2} + \frac{\eu^x -
\eu^{-x}}{2} = \eu^x$. Portanto, para $n \in \Z$:
\[
(\cosh x + \sinh x)^n = (\eu^x)^n = \eu^{nx}
= \cosh nx + \sinh nx .
\]
\end{solution}

\begin{solution}{exo:b1:functions:4}
Com $y = \arcsin x$: $\cos 2y = 1 - 2\sin^2 y = 1 - 2x^2$, logo
$\cos(2\arcsin x) = 1 - 2x^2$.

Com $y = \arctan x$: $\sin 2y = 2 \sin y \cos y = 2 \tan y \cos^2 y =
\frac{2\tan y}{1 + \tan^2 y}$, logo $\sin(2\arctan x) =
\dfrac{2x}{1 + x^2}$.
\end{solution}

\begin{solution}{exo:b1:functions:5}
Do crescimento mais lento ao mais rápido:
\[
(\ln x)^{1000} \;\ll\; x^{100} \;\ll\; x^{\ln x} \;\ll\;
\eu^{\sqrt x} \;\ll\; \eu^{x/100} .
\]
Justificativas: $(\ln x)^{1000}/x^{100} \to 0$ pela
\cref{prop:b1:functions:powerrules} (os logaritmos perdem para as potências).
$x^{100} = \eu^{100\ln x}$ e $x^{\ln x} = \eu^{(\ln x)^2}$: como
$(\ln x)^2 - 100\ln x \to +\infty$, a segunda vence. $x^{\ln x} =
\eu^{(\ln x)^2} \ll \eu^{\sqrt x}$ porque $(\ln x)^2/\sqrt x \to 0$
(os logaritmos perdem para a potência $x^{1/4}$, elevada ao quadrado).
Por fim, $\sqrt x - x/100
\to -\infty$, de modo que $\eu^{\sqrt x} \ll \eu^{x/100}$.
\end{solution}

\begin{solution}{exo:b1:functions:6}
Domínio: $x \neq \pm 1$, três intervalos. Em cada um deles,
\[
f'(x) = \frac{\bigl(\frac{2x}{1-x^2}\bigr)'}{1 +
\frac{4x^2}{(1-x^2)^2}}
= \frac{\frac{2(1-x^2) + 4x^2}{(1-x^2)^2}}
{\frac{(1-x^2)^2 + 4x^2}{(1-x^2)^2}}
= \frac{2(1 + x^2)}{(1 + x^2)^2} = \frac{2}{1 + x^2}
= (2\arctan x)' .
\]
Logo, $f(x) - 2\arctan x$ é constante em cada intervalo. Valores: em $x =
0$, $f(0) = 0$: a constante é $0$ em $\intoo{-1}{1}$. Quando $x \to
+\infty$, $\frac{2x}{1-x^2} \to 0^-$, de modo que $f \to 0$, ao passo que $2\arctan x
\to \pi$: a constante é $-\pi$ em $\intoo{1}{+\infty}$. Pela imparidade,
ela é $+\pi$ em $\intoo{-\infty}{-1}$. Em resumo: $f = 2\arctan x$ em
$\intoo{-1}{1}$, $f = 2\arctan x - \pi$ para $x > 1$, $f = 2\arctan x +
\pi$ para $x < -1$. (Trata-se da fórmula do arco duplo da tangente,
lida através de $\arctan$.)
\end{solution}

\begin{solution}{exo:b1:functions:7}
$\cosh x = 2$: com $u = \eu^x > 0$, $u + u^{-1} = 4$, de modo que $u^2 - 4u +
1 = 0$, $u = 2 \pm \sqrt 3$: $x = \ln(2 + \sqrt 3)$ ou $x = \ln(2 -
\sqrt 3) = -\ln(2 + \sqrt 3)$ (as duas soluções são opostas, pois
$\cosh$ é par; ambas são válidas). Equivalentemente, $x = \pm
\operatorname{arcosh} 2$.

$5\cosh x - 4\sinh x = 3$: substituindo as definições exponenciais,
$\frac{5(u + u^{-1}) - 4(u - u^{-1})}{2} = 3$, isto é, $u + 9u^{-1} =
6$, ou seja, $u^2 - 6u + 9 = (u - 3)^2 = 0$: $u = 3$, uma única solução $x =
\ln 3$.
\end{solution}

\begin{solution}{exo:b1:functions:8}
Seja $\alpha = \arctan\frac12 + \arctan\frac13$. Fórmula de adição:
\[
\tan\alpha = \frac{\frac12 + \frac13}{1 - \frac12\cdot\frac13}
= \frac{5/6}{5/6} = 1 .
\]
Os dois \omterm{def:b1:functions:arc}{arco-tangentes} estão em $\intoo{0}{\frac\pi4}$ (seus argumentos estão em
$\intoo{0}{1}$), de modo que $\alpha \in \intoo{0}{\frac\pi2}$; o único ângulo
ali com tangente $1$ é $\frac\pi4$.

Machin: seja $\beta = \arctan\frac15$. Arco duplo duas vezes:
\[
\tan 2\beta = \frac{2/5}{1 - 1/25} = \frac{5}{12},
\qquad
\tan 4\beta = \frac{2 \cdot 5/12}{1 - 25/144} = \frac{120}{119}.
\]
Então
\[
\tan\Bigl(4\beta - \frac\pi4\Bigr)
= \frac{\frac{120}{119} - 1}{1 + \frac{120}{119}}
= \frac{1/119}{239/119} = \frac{1}{239}.
\]
Localização: $\beta < \arctan 1 = \frac\pi4$ e, de fato, $\tan 4\beta =
\frac{120}{119}$ está próximo de $1$, com $4\beta \in \intoo{0}{\frac\pi2}$
(pois $\beta < \frac\pi8$, já que $\tan\frac\pi8 = \sqrt 2 - 1 > \frac15$);
logo $4\beta - \frac\pi4 \in \intoo{-\frac\pi4}{\frac\pi4}$, onde
$\arctan$ inverte $\tan$: $4\beta - \frac\pi4 =
\arctan\frac{1}{239}$, que é a fórmula de Machin.
\end{solution}

\begin{solution}{exo:b1:functions:9}
Seja $f(x) = x - \sin x$: $f(0) = 0$ e $f'(x) = 1 - \cos x \geq 0$,
de modo que $f \geq 0$ em $\R_+$: $\sin x \leq x$.

Seja $g(x) = \sin x - x + \frac{x^3}{6}$: $g(0) = 0$, $g'(x) = \cos x -
1 + \frac{x^2}{2}$, $g'(0) = 0$, $g''(x) = -\sin x + x = f(x) \geq 0$
em $\R_+$. Assim, $g'$ é crescente com $g'(0) = 0$, logo $g' \geq 0$,
logo $g$ é crescente com $g(0) = 0$: $g \geq 0$, isto é, $x -
\frac{x^3}{6} \leq \sin x$ em $\R_+$.

Consequentemente, $0 \leq \frac{x - \sin x}{x^3} \leq \frac 16$ para $x >
0$: a razão fica presa na escala de $\frac16$, e o
\cref{ch:b1:taylor} mostra que o seu limite é exatamente $\frac 16$.
\end{solution}

\begin{solution}{exo:b1:functions:10}
Escreva $y = \arcsin x$, de modo que $2x\sqrt{1 - x^2} = 2\sin y\cos y = \sin
2y$ e $h(x) = \arcsin(\sin 2y) - 2y$.

Para $\abs x \leq \frac{\sqrt 2}{2}$: $2y \in
\intcc{-\frac\pi2}{\frac\pi2}$, de modo que $\arcsin(\sin 2y) = 2y$ e $h =
0$.

Para $x \in \intoc{\frac{\sqrt 2}{2}}{1}$: $2y \in
\intoc{\frac\pi2}{\pi}$, e o ângulo de
$\intcc{-\frac\pi2}{\frac\pi2}$ cujo seno vale $\sin 2y$ é $\pi - 2y$:
logo $h(x) = \pi - 4\arcsin x$. (Verificação pela derivada: ali
$h'(x) = \frac{2(1 - 2x^2)}{\abs{1 - 2x^2}\sqrt{1 - x^2}} -
\frac{2}{\sqrt{1-x^2}} = \frac{-4}{\sqrt{1 - x^2}}$, a derivada
de $-4\arcsin x$; e em $x = 1$, $h(1) = \arcsin 0 - 2\cdot\frac\pi2
= -\pi = \pi - 4\cdot\frac\pi2$.)

Para $x \in \intco{-1}{-\frac{\sqrt 2}{2}}$, pela imparidade de $h$:
$h(x) = -\pi - 4\arcsin x$.
\end{solution}

\begin{solution}{exo:b1:functions:11}
$g = \arctan \circ \sinh$ é uma composta de funções ímpares e estritamente
crescentes, logo é ímpar e estritamente crescente; quando $x \to
+\infty$, $\sinh x \to +\infty$, de modo que $g(x) \to \frac\pi2$, e $g$ é uma
bijeção sobre $\intoo{-\frac\pi2}{\frac\pi2}$ (continuidade mais o
teorema do valor intermediário). Regra da cadeia com a
\cref{prop:b1:functions:hyprules}~(1):
\[
g'(x) = \frac{\cosh x}{1 + \sinh^2 x} = \frac{\cosh x}{\cosh^2 x}
= \frac{1}{\cosh x} .
\]
Por construção, $\tan g(x) = \sinh x$. Então, como $g(x) \in
\intoo{-\frac\pi2}{\frac\pi2}$ tem cosseno positivo,
\begin{align*}
\cos g(x) &= \frac{1}{\sqrt{1 + \tan^2 g(x)}}
= \frac{1}{\sqrt{1 + \sinh^2 x}} = \frac{1}{\cosh x},\\
\sin g(x) &= \tan g(x) \cos g(x) = \frac{\sinh x}{\cosh x} = \tanh x .
\end{align*}
\end{solution}

\begin{solution}{exo:b1:functions:12}
Ponha $u = \arctan 2$ e $v = \arctan 3$; ambos estão em
$\intoo{\frac\pi4}{\frac\pi2}$ (seus argumentos excedem $1$), de modo que $u
+ v \in \intoo{\frac\pi2}{\pi}$. A fórmula de adição dá
\[
\tan(u + v) = \frac{2 + 3}{1 - 6} = -1 ,
\]
e o único ângulo de $\intoo{\frac\pi2}{\pi}$ com tangente $-1$
é $\frac{3\pi}4$: logo $\arctan 2 + \arctan 3 = \frac{3\pi}4$.
(Aplicar $\arctan$ cegamente à tangente daria
$-\frac\pi4$, errado por $\pi$ --- a localização é o que salva o
cálculo, e o denominador negativo $1 - 6$ é precisamente o
sinal de que a soma saiu de
$\intoo{-\frac\pi2}{\frac\pi2}$.) Somando $\arctan 1 = \frac\pi4$:
\[
\arctan 1 + \arctan 2 + \arctan 3 = \frac\pi4 + \frac{3\pi}4 = \pi .
\]
\end{solution}

\begin{solution}{pb:b1:functions:1}
\textbf{1.} $f$ é um quociente de funções deriváveis com
denominador não nulo em $\intoo0{+\infty}$, e
\[
f'(t) = \frac{\frac1t \cdot t - \ln t}{t^2} = \frac{1 - \ln
t}{t^2} ,
\]
positiva para $t < \eu$, nula em $t = \eu$, negativa para $t > \eu$:
$f$ cresce estritamente em $\intoc0\eu$, decresce estritamente em
$\intco\eu{+\infty}$, com máximo $f(\eu) = \frac1\eu$.

\textbf{2.} Quando $t \to 0^+$: $\ln t \to -\infty$ e $\frac1t \to
+\infty$, de modo que $f(t) \to -\infty$. Quando $t \to +\infty$: $f(t) \to 0$
pela \omterm{prop:b1:functions:powerrules}{comparação de crescimentos} da \cref{prop:b1:functions:powerrules}
($\beta = \alpha = 1$). Sinal: o de $\ln t$, de modo que $f < 0$ em
$\intoo01$, $f(1) = 0$, $f > 0$ em $\intoo1{+\infty}$. O gráfico
sobe de $-\infty$, cruza o zero em $1$, atinge o pico em $(\eu,
\frac1\eu)$ e depois decai a $0^+$.

\textbf{3.} $f(4) = \frac{\ln 4}4 = \frac{2\ln 2}4 = \frac{\ln 2}2
= f(2)$.

\textbf{4.} Pela tabela de variação e pelo teorema do valor
intermediário (usado no nível do ensino médio; formalizado no
\cref{ch:b1:continuity}). Para $c < 0$: só há soluções onde
$f < 0$, isto é, em $\intoo01$, onde $f$ é uma bijeção estritamente
crescente sobre $\intoo{-\infty}0$: exatamente uma. Para $c = 0$: apenas
$t = 1$. Para $0 < c < \frac1\eu$: em $\intoo1\eu$, $f$ cresce
de $0$ até $\frac1\eu$: uma solução; em $\intoo\eu{+\infty}$,
$f$ decresce de $\frac1\eu$ a $0$: mais uma; total, duas. Para
$c = \frac1\eu$: apenas o ponto de máximo $t = \eu$. Para $c >
\frac1\eu$: nenhuma.

\textbf{5.} Para $x \neq \eu$, a maximalidade estrita dá $f(x) <
f(\eu)$, isto é, $\frac{\ln x}x < \frac1\eu$, ou seja, $\eu\ln x < x$,
isto é, $\ln(x^\eu) < \ln(\eu^x)$: $x^\eu < \eu^x$. Com $x = \pi$:
$\eu^\pi > \pi^\eu$ (numericamente, $23.14 > 22.46$).

\textbf{6.} Para $x, y > 0$, os dois lados são positivos, de modo que
\[
x^y = y^x \iff y\ln x = x\ln y \iff \frac{\ln x}x = \frac{\ln y}y
\iff f(x) = f(y) ,
\]
dividindo por $xy > 0$.

\textbf{7.} Sejam $f(x) = f(y) = c$ com $x \neq y$. Se $c \leq 0$ ou
$c = \frac1\eu$, a questão 4 diz que a equação $f = c$ tem uma única
solução: impossível. Logo $0 < c < \frac1\eu$ e, de novo pela
questão 4, as duas soluções são um ponto de $\intoo1\eu$ e um
de $\intoo\eu{+\infty}$: as duas coordenadas excedem $1$ e elas
ficam de lados opostos de $\eu$.

\textbf{8.} Substituindo $y = tx$ em $y\ln x = x\ln y$:
\[
tx\ln x = x(\ln t + \ln x)
\iff (t - 1)\ln x = \ln t
\iff \ln x = \frac{\ln t}{t - 1} ,
\]
de modo que $x = t^{1/(t-1)}$ e $y = tx = t^{1 + \frac1{t-1}} =
t^{t/(t-1)}$. Reciprocamente, para esses valores, $\ln y = \frac{t\ln
t}{t-1} = t\ln x$ e $y = tx$, logo $\frac{\ln y}y = \frac{t\ln
x}{tx} = \frac{\ln x}x$: uma solução. Toda solução não trivial tem
alguma razão $t = y/x \in \intoo0{+\infty} \setminus \{1\}$, de modo que a
parametrização é completa.

\textbf{9.} $t = 2$: $x = 2^{1/1} = 2$, $y = 2^{2/1} = 4$. E
\[
x(1/t) = (1/t)^{\frac1{\frac1t - 1}}
= (t^{-1})^{\frac{t}{1 - t}}
= t^{\frac{t}{t - 1}} = y(t) ,
\]
e então $y(1/t) = \frac1t\,x(1/t) = \frac{y(t)}t = x(t)$: a
mudança de parâmetro $t \mapsto 1/t$ troca as coordenadas, como a
simetria da equação exige.

\textbf{10.} Quando $t \to 1$: $\frac{\ln t}{t-1} \to 1$ (é o
quociente de diferenças de $\ln$ em $1$), de modo que $x \to \eu^1 = \eu$ e
$y = tx \to \eu$. Quando $t \to +\infty$: $\ln x = \frac{\ln t}{t-1}
\to 0$, logo $x \to 1$, ao passo que $\ln y = \frac{t\ln t}{t-1} \to
+\infty$, logo $y \to +\infty$. Quando $t \to 0^+$: $\ln x = \frac{\ln
t}{t-1} \to \frac{-\infty}{-1} = +\infty$, logo $x \to +\infty$,
ao passo que $\ln y = \frac{t\ln t}{t-1} \to \frac{0}{-1} = 0$, logo $y \to
1$ (usando $t\ln t \to 0$, \cref{ex:b1:functions:xx}). O ramo,
portanto, vai da assíntota $y = 1$ (extrema direita), sobe passando por
$(\eu, \eu)$ na diagonal e escapa ao longo da assíntota $x = 1$
(extremo superior) --- simétrico em relação à diagonal pela questão 9.

\textbf{11.} $\ln x(t) = g(t) = \frac{\ln t}{t-1}$, e
\[
g'(t) = \frac{\frac{t-1}t - \ln t}{(t-1)^2}
= \frac{h(t)}{(t-1)^2},
\qquad h(t) = 1 - \frac1t - \ln t .
\]
$h(1) = 0$ e $h'(t) = \frac1{t^2} - \frac1t = \frac{1 -
t}{t^2} < 0$ para $t > 1$: $h < 0$ em $\intoo1{+\infty}$, de modo que $g' <
0$ e $x = \eu^g$ é estritamente decrescente ali (de $\eu$ até
$1$, pela questão 10).

\textbf{12.} Pela questão 11, $t \mapsto x(t)$ é uma bijeção
estritamente decrescente de $\intoo1{+\infty}$ sobre $\intoo1\eu$;
escreva $t(x)$ para a sua inversa (também estritamente decrescente) e
ponha $\varphi(x) = y(t(x))$. Note que $h < 0$ também em $\intoo01$ (ali
$h' > 0$ e $h(1) = 0$), de modo que $x(\cdot)$ é decrescente em
$\intoo01$ também; portanto, $y(t) = x(1/t)$ (questão 9) é
\emph{crescente} em $t$ em $\intoo1{+\infty}$, de $\eu$ a
$+\infty$. Compondo: $\varphi = y \circ t(\cdot)$ é estritamente
decrescente de $\intoo1\eu$ sobre $\intoo\eu{+\infty}$, e
$x^{\varphi(x)} = \varphi(x)^x$ pela questão 8. Por fim, para um
valor comum $c = f(x) \in \intoo0{\frac1\eu}$, o par $\{x,
\varphi(x)\}$ é \emph{o} \omterm{def:b1:logic:sets}{conjunto} solução de dois elementos de $f = c$
(questão 4); estendendo $\varphi$ a $\intoo\eu{+\infty}$ como a
\omterm{def:b1:logic:map}{aplicação} inversa, $\varphi(\varphi(x))$ devolve o outro elemento (isto é,
o original): $\varphi \circ \varphi = \mathrm{id}$.

\textbf{13.} Se $(x, y)$ é uma solução inteira não trivial com $x
< y$, a questão 7 põe $x \in \intoo1\eu$: o único inteiro ali
é $x = 2$. Então $f(y) = f(2) = \frac{\ln2}2$ com $y > \eu$; pela
questão 4, a equação $f = \frac{\ln2}2$ tem exatamente uma
solução além de $\eu$, e a questão 3 a exibe: $y = 4$. Portanto,
$(2, 4)$ e a sua troca são as únicas.

\textbf{14.} $t = 1 + \frac1n = \frac{n+1}n$ dá $\frac1{t - 1}
= n$ e $\frac t{t-1} = n + 1$, de modo que
\[
x_n = \Bigl(\frac{n+1}n\Bigr)^{n},
\qquad
y_n = \Bigl(\frac{n+1}n\Bigr)^{n+1},
\]
manifestamente racionais, distintos ($y_n = t\,x_n \neq x_n$), e $n =
1$ dá $x_1 = 2$, $y_1 = 4$.

\textbf{15.} $t = y/x$ é racional e $> 1$; escreva $t - 1 =
\frac rs$ na forma irredutível, de modo que $t = \frac{s + r}s$ com $\gcd(s +
r, s) = \gcd(r, s) = 1$. A questão 8 dá $x = t^{s/r}$, logo
\[
x^r = t^s = \frac{(s+r)^s}{s^s} ,
\]
uma fração irredutível (nenhum primo divide simultaneamente $s + r$ e
$s$). Escrevendo $x = \frac pq$ na forma irredutível, $x^r = \frac{p^r}
{q^r}$ também é irredutível e, pela unicidade da
representação reduzida: $p^r = (s+r)^s$ e $q^r = s^s$. Compare o
expoente de um primo qualquer $\ell$ em $q^r = s^s$: $r \cdot
v_\ell(q) = s \cdot v_\ell(s)$, de modo que $r$ divide $s\,v_\ell(s)$;
como $\gcd(r, s) = 1$, $r$ divide $v_\ell(s)$ para todo primo
$\ell$ (fatoração única, admitida; demonstrada no
\cref{ch:b1:arith}), logo $s = b^r$ para um inteiro $b$. O mesmo
argumento em $p^r = (s+r)^s$ dá $s + r = a^r$.

\textbf{16.} Suponha $a^r - b^r = r$ com inteiros $a > b \geq 1$
e $r \geq 2$. Então $a \geq b + 1$, de modo que, pelo teorema binomial,
\[
a^r - b^r \geq (b+1)^r - b^r
= \sum_{k=0}^{r-1}\binom rk b^k
\geq \binom r0 + \binom r{r-1} b^{r-1}
\geq 1 + r > r ,
\]
uma contradição. Da questão 15, $a^r - b^r = (s + r) - s = r$
força $r = 1$: $t = 1 + \frac1s$, e a solução é $(x_s,
y_s)$ da questão 14. Junto com os pares trocados, a
classificação está completa.

\textbf{17.} $f\bigl(\tfrac94\bigr) = \frac{\ln 2.25}{2.25} =
\frac{0.81093}{2.25} = 0.3604$ e $f\bigl(\tfrac{27}8\bigr) =
\frac{\ln 3.375}{3.375} = \frac{1.21640}{3.375} = 0.3604$ (quatro
casas decimais): iguais, como a construção promete --- de modo que
$\bigl(\tfrac94\bigr)^{27/8} = \bigl(\tfrac{27}8\bigr)^{9/4}$.

\textbf{18.} Os parâmetros $t_n = 1 + \frac1n$ decrescem estritamente
para $1$. Como $x(\cdot)$ é estritamente decrescente em
$\intoo1{+\infty}$ (questão 11), $x_n = x(t_n)$ é estritamente
\emph{crescente}; como $y(\cdot)$ é estritamente crescente ali
(questão 12), $y_n = y(t_n)$ é estritamente \emph{decrescente}. Pela
questão 7, $x_n \in \intoo1\eu$ e $y_n \in
\intoo\eu{+\infty}$: logo $x_n < \eu < y_n$ para todo $n$. Por fim,
$t_n \to 1$ e a questão 10 dão $x_n \to \eu$ e $y_n \to \eu$.
A clássica convergência monótona de $\bigl(1 +
\frac1n\bigr)^n$ cai da geometria do ramo, sem nenhuma
desigualdade nova.

\textbf{19.} Se $\eu \leq a < b$: $f$ estritamente decrescente em
$\intco\eu{+\infty}$ dá $f(a) > f(b)$, isto é, $b\ln a > a\ln
b$, ou seja, $a^b > b^a$. Se $1 < a < b \leq \eu$: $f$ estritamente
crescente dá $f(a) < f(b)$, logo $a^b < b^a$. Caso misto:
$2 < \eu < 3$ e $2^3 = 8 < 9 = 3^2$; mas $2 < \eu < 5$ e
$2^5 = 32 > 25 = 5^2$: os dois resultados ocorrem, decididos por qual lado
do ramo o ponto $(a, b)$ cai.

\textbf{20.} O ramo encontra a diagonal em $\eu$, e $\varphi$
se estende continuamente ali com $\varphi(\eu) = \eu$.
Derivando $\varphi(\varphi(x)) = x$ em $x = \eu$ pela
regra da cadeia: $\varphi'(\varphi(\eu))\,\varphi'(\eu) =
\varphi'(\eu)^2 = 1$, de modo que $\varphi'(\eu) = \pm1$; $\varphi$ é
decrescente, logo $\varphi'(\eu) = -1$. O ramo cruza a
diagonal com coeficiente angular $-1$: perpendicularmente.

\textbf{21.} $y = 3x$ é o caso $t = 3$: $x = 3^{1/2} = \sqrt3$
e $y = 3^{3/2} = 3\sqrt3$. Verificação: $f(\sqrt3) =
\frac{0.54931}{1.73205} = 0.3171$ e $f(3\sqrt3) =
\frac{\ln 5.19615}{5.19615} = \frac{1.64792}{5.19615} = 0.3171$:
iguais até a quarta casa decimal, de modo que $(\sqrt3)^{3\sqrt3} =
(3\sqrt3)^{\sqrt3}$.

\textbf{22.} $n^{1/n} = \eu^{f(n)}$, e $\exp$ é crescente, de modo que
a ordenação é a de $f(n)$: $f(3) = \frac{\ln3}3 \approx
0.3662$ supera $f(2) = f(4) = \frac{\ln2}2 \approx 0.3466$ (questão
3) e $f(5) \approx 0.3219$. O maior é $\sqrt[3]3$, e o
par de iguais é $\sqrt2 = \sqrt[4]4$ --- o par inteiro $(2, 4)$
em mais um disfarce.

\textbf{23.} O \omterm{def:b1:logic:sets}{conjunto} solução é a união da diagonal $\{(x,
x) : x > 0\}$ com um único ramo, simétrico em relação à diagonal,
estritamente decrescente da assíntota $x = 1$ (quando $y \to
+\infty$) até a assíntota $y = 1$ (quando $x \to +\infty$), cruzando
a diagonal exatamente uma vez, em $(\eu, \eu)$, ali com coeficiente
angular $-1$. Sobre o ramo estão os pontos inteiros $(2, 4)$ e $(4, 2)$
--- os únicos --- e os pontos racionais $(x_n, y_n)$,
que marcham monotonamente ao longo do ramo rumo a $(\eu, \eu)$
sem jamais alcançá-lo ($\eu$ é irracional; os pontos racionais do ramo
se acumulam num canto irracional).

\textbf{24.} (i) As comparações de crescimento deram os limites de $f$ em
$+\infty$ (questão 2) e $t\ln t \to 0$ (questão 10),
moldando tanto a tabela de variação quanto as assíntotas. (ii) O
teorema do valor intermediário, através dos enunciados de bijeção,
converteu a tabela de variação em contagens exatas de soluções
(questão 4) e na existência da \omterm{def:b1:logic:map}{aplicação} inversa $t(x)$
(questão 12). (iii) A fatoração única alimentou as duas
identificações de frações irredutíveis da questão 15, o coração
aritmético da classificação dos pontos racionais.

\textbf{25.} Um único cálculo de derivada --- o sinal de $1 - \ln
t$ --- gerou tudo: a contagem de soluções para cada nível
$c$, a forma e as assíntotas do ramo, a desigualdade
extremal $\eu^x \geq x^\eu$, a classificação das soluções inteiras e
racionais, e a convergência monótona de $\bigl(1 +
\frac1n\bigr)^n$. Esta é a economia de pensar com tabelas de
variação: um estudo unidimensional resolve uma equação em duas
variáveis porque a equação se fatora por uma única função.
O \cref{ch:b1:seq} refará a convergência de $(x_n)$ com a
teoria geral das sequências monótonas; o \cref{ch:b1:derivative}
funda rigorosamente as tabelas de variação no teorema do valor médio; e
o \cref{ch:b1:taylor} fornece o que a tabela não pode --- a
\emph{velocidade} da convergência $x_n \to \eu$ (ela é de ordem
$1/n$).
\end{solution}
